\documentclass[11pt]{article}

%%%%%%%%%%%%%%Include Packages%%%%%%%%%%%%%%%%%%%%%%%%%%
\usepackage{xcolor}
\usepackage{mathtools}
\usepackage[a4paper, total={6in, 8in}, margin=1.25in]{geometry}
\usepackage{amsmath}
\usepackage{amssymb}
\usepackage{paralist}
\usepackage{rsfso}
\usepackage{amsthm}
\usepackage{wasysym}
\usepackage[inline]{enumitem}   
\usepackage{hyperref}
\usepackage{tocloft}
\usepackage{titlesec}
\usepackage{colortbl}
\usepackage{wrapfig}
\usepackage{pdfpages}
%%%%%%%%%%%%%%%%%%%%%%%%%%%%%%%%%%%%%%%%%%%%%%%%%%%%%%%%



\begin{document}

\begin{wrapfigure}[8]{r}{0.35\textwidth}
  \begin{center}
  \vspace{-15pt}
    \includegraphics[scale=0.6]{pumpkin1}
  \end{center}
\end{wrapfigure}
\Large \noindent \textbf{Physics 5A Discussion (Week 5)}\\
\normalsize
Department of Physics and Astronomy, UCLA\\
\noindent\rule{8cm}{0.4pt}\\

\noindent (a) Let's consider a $2\,$kg pumpkin in free-fall motion. We set up the $x$- and $y$-axes as shown. Find the \textit{net force} (the sum of all forces, but we only have the gravitational force in this case) acting on the pumpkin, write it in vector-component form.
\hfill\newline\hfill\newline
\hfill\newline\hfill\newline
\hfill\newline\hfill\newline



\begin{wrapfigure}[8]{r}{0.4\textwidth}
  \begin{center}
    \includegraphics[scale=0.6]{pumpkin2}
  \end{center}
\end{wrapfigure}
\hfill\break
\hfill\break
\hfill\break
\noindent (b) Now the pumpkin sits steadily on the table.
Find the net force acting on the pumpkin (which includes both the gravitational and normal forces in this case), write it in vector-component form.
\hfill\newline\hfill\newline
\hfill\newline\hfill\newline
\hfill\newline\hfill\newline
\hfill\newline\hfill\newline
\hfill\newline\hfill\newline
\hfill\newline\hfill\newline


\noindent (c) Draw the free-body diagram of the pumpkin on the table. There are only two forces acting on the pumpkin: the gravitational force by gravity, and the normal force exerted by the table. Since the pumpkin is not moving, these two forces should cancel out.
\hfill\newline\hfill\newline
\hfill\newline\hfill\newline
\hfill\newline\hfill\newline

\newpage
\noindent Now let's consider the same $2\,$kg pumpkin, sliding down a tilted table. Use a new coordinate system in which the $x$- and $y$-axes are tilted at an angle $\theta=30^\circ$ as shown. \\

\begin{center}
\includegraphics[scale=0.6]{pumpkin3}
\end{center}

\noindent (c) Decompose the gravitational force acting on the pumpkin into components as shown, and write it in component form. (\textit{Hint: Identify the angle in the triangle, then find the $\perp$-component and $\parallel$-component using trigonometry.})
\hfill\newline\hfill\newline
\hfill\newline\hfill\newline
\hfill\newline\hfill\newline
\hfill\newline\hfill\newline

\noindent (d) Complete the free-body diagram of the pumpkin. The gravitational force has been drawn for you. Add the normal force exerted by the table, and be careful about its magnitude and direction.
\hfill\newline\hfill\newline
\hfill\newline\hfill\newline
\hfill\newline\hfill\newline
\hfill\newline\hfill\newline

\noindent (e) The pumpkin slides down the table. Ignoring any friction, find the net acceleration of the pumpkin and write it in component form in the tilted coordinate system. (\textit{Hint: The $?$-component of the gravitational force gets canceled, so we only have the $?$-component of the gravitational force left.})


\includepdf[page=-]{wk5s}




\end{document}


